% Fichier généré automatiquement par tools/generate_corrections.py.
% Source : CIN/CIN-02-VitesseAcceleration/09_RT_RSG/09_RT_RSG.tex
% Statut : complété (4/4 question(s)).
\begin{corrigebox}[Corrigé complété]
\CorrigeQuestion{1}
$\vectv{B}{2}{0}=\vectv{B}{2}{1}+\vectv{B}{1}{0}$.

D'une part,  $\vectv{B}{2}{1}=\lambdap\vi{1}$.

D'autre part, en utilisant le roulement sans glissement en $I$, 
$\babarv{B}{I}{1}{0} $ $=\vect{0}+\left(-\lambda(t)\vi{1}-R\vj{0} \right)\wedge\thetap \vk{0}$
$=-\thetap\left(\lambda(t)\vi{1}\wedge \vk{0}+R\vj{0}\wedge \vk{0} \right)$
$=\thetap\left(\lambda(t)\vj{1}-R\vi{0} \right)$.

Au final, $\vectv{B}{2}{0} = \lambdap\vi{1} + \thetap\left(\lambda(t)\vj{1}-R\vi{0} \right)$.
\CorrigeQuestion{2}
On construit un sommet par solide, puis une arête par liaison en
        indiquant son type, son centre et son axe. Les actions extérieures (poids,
        actionneur, charge) sont reliées au solide concerné. Le graphe doit retrouver le
        nombre de mobilités et de boucles du mécanisme de l'énoncé.
\CorrigeQuestion{3}
$\torseurcin{V}{2}{0}=\torseurl{\thetap \vk{0}}{ \lambdap\vi{1} + \thetap\left(\lambda(t)\vj{1}-R\vi{0} \right)}{B}$.
\CorrigeQuestion{4}
$\vectg{B}{2}{0} = \deriv{\vectv{B}{2}{0}}{\rep{0}}$
$ = \lambdapp(t)\vi{1} +\lambdap(t)\thetap\vj{1} 
+ \thetapp(t)\left(\lambda(t)\vj{1}-R\vi{0} \right)
+ \thetap(t)\left(\lambdap(t)\vj{1}-\lambda(t)\thetap\vi{1} \right)
$.

\begin{enumerate}
\item $\vectv{B}{2}{0} = \lambdap\vi{1} + \thetap\left(\lambda(t)\vj{1}-R\vi{0} \right)$.
\item $\torseurcin{V}{2}{0}=\torseurl{\thetap \vk{0}}{ \lambdap\vi{1} + \thetap\left(\lambda(t)\vj{1}-R\vi{0} \right)}{B}$.
\item $\vectg{B}{2}{0}  = \lambdapp(t)\vi{1} +\lambdap(t)\thetap\vj{1} 
+ \thetapp(t)\left(\lambda(t)\vj{1}-R\vi{0} \right)
+ \thetap(t)\left(\lambdap(t)\vj{1}-\lambda(t)\thetap\vi{1} \right)
$.
\end{enumerate}
\end{corrigebox}
